% GENERATED FILE. Edit canonical lessons, then run npm run generate:books.
% canonical-source-hash: fnv1a64:a7b189644455d18e
% canonical-lessons: GE-C446-reinigung, GE-C446-waschsalon, GE-C446-entfernen, GE-C446-flicken, GE-C446-trocknen

\chapter{In der Reinigung}
\label{ch:ge-in-der-reinigung}
\hlchaptermodality{446}

\begin{chapteropening}
\textbf{By the end of this chapter:} \emph{I can get clothes cleaned and altered: the dry cleaner, the launderette, removing a stain, patching and drying.}

Complete the last lesson of chapter 446: I can get clothes cleaned and altered: the dry cleaner, the launderette, removing a stain, patching and drying.
\end{chapteropening}

\section[die Reinigung]{die Reinigung — the dry cleaner; the cleaning}

\label{lesson:GE-C446-reinigung}



\begin{quote}
Before the new one: say the German for the locksmith, then the German for to unlock.
\end{quote}

\subsection*{You'll want to know: die Reinigung}
\textbf{die Reinigung} — \textquotedblleft{}the dry cleaner; the cleaning\textquotedblright{}.

\begin{quote}
Den Anzug bringt er in die Reinigung.
\end{quote}

\begin{quote}
\emph{He takes the suit to the dry cleaner.}
\end{quote}

Plural: \emph{die Reinigungen}. Also the cleaning of a room.

\begin{cousinweb}
\emph{Reinigung} is built on \emph{reinigen} 'to clean', from \emph{rein} 'clean, pure', inherited from Proto-Germanic \emph{hrainiz} (EtymDB lists \emph{Reinigung}).
\end{cousinweb}

\begin{grammarlens}[title={What you've built}]
One more word for telling, explaining and describing at B1.
\end{grammarlens}

\subsection*{Your turn}
\begin{itemize}
\raggedright
  \item \emph{Say it:} \emph{die Reinigung}
  \item \emph{Say it:} \emph{die Reinigung}, once more
  \item \emph{Say it:} a sentence of your own with \emph{die Reinigung}
  \item \emph{Recall:} say the German for the locksmith, then the German for to unlock, then say \emph{die Reinigung} again
\end{itemize}

\begin{tcolorbox}[breakable,colback=teal!4,colframe=teal!35!black,title={Before you move on}]
What does die Reinigung mean? (\textquotedblleft{}The dry cleaner; the cleaning\textquotedblright{}.) Say it once more.
\end{tcolorbox}
\section[der Waschsalon]{der Waschsalon — the launderette (US: laundromat)}

\label{lesson:GE-C446-waschsalon}



\begin{quote}
Before the new one: say the German for to unlock, then the German for the dry cleaner.
\end{quote}

\subsection*{You'll want to know: der Waschsalon}
\textbf{der Waschsalon} — \textquotedblleft{}the launderette (US: laundromat)\textquotedblright{}.

\begin{quote}
Im Waschsalon liest sie ein Buch, während die Maschine läuft.
\end{quote}

\begin{quote}
\emph{In the launderette she reads a book while the machine runs.}
\end{quote}

Plural: \emph{die Waschsalons}.

\begin{cousinweb}
\emph{Waschsalon} is built from \emph{waschen} 'to wash', inherited from Proto-Germanic \emph{waskaną}, and \emph{Salon} 'room', borrowed from French \emph{salon}, which came from Italian \emph{salone} 'large hall' (EtymDB).
\end{cousinweb}

\begin{grammarlens}[title={What you've built}]
One more word for telling, explaining and describing at B1.
\end{grammarlens}

\subsection*{Your turn}
\begin{itemize}
\raggedright
  \item \emph{Say it:} \emph{der Waschsalon}
  \item \emph{Say it:} \emph{der Waschsalon}, once more
  \item \emph{Say it:} a sentence of your own with \emph{der Waschsalon}
  \item \emph{Recall:} say the German for to unlock, then the German for the dry cleaner, then say \emph{der Waschsalon} again
\end{itemize}

\begin{tcolorbox}[breakable,colback=teal!4,colframe=teal!35!black,title={Before you move on}]
What does der Waschsalon mean? (\textquotedblleft{}The launderette (US: laundromat)\textquotedblright{}.) Say it once more.
\end{tcolorbox}
\section[entfernen]{entfernen — to remove}

\label{lesson:GE-C446-entfernen}



\begin{quote}
Before the new one: say the German for the dry cleaner, then the German for the launderette.
\end{quote}

\subsection*{You'll want to know: entfernen}
\textbf{entfernen} — \textquotedblleft{}to remove\textquotedblright{}.

\begin{quote}
Die Reinigung entfernt den Fleck vom Kleid.
\end{quote}

\begin{quote}
\emph{The dry cleaner removes the stain from the dress.}
\end{quote}

Formal. \emph{Entfernt} also means distant: \emph{zwei Kilometer entfernt}.

\begin{cousinweb}
\emph{entfernen} is inherited from Middle High German \emph{entverren} (EtymDB), built on \emph{fern} 'far': to put something far away.
\end{cousinweb}

\begin{grammarlens}[title={What you've built}]
One more word for telling, explaining and describing at B1.
\end{grammarlens}

\subsection*{Your turn}
\begin{itemize}
\raggedright
  \item \emph{Say it:} \emph{entfernen}
  \item \emph{Say it:} \emph{entfernen}, once more
  \item \emph{Say it:} a sentence of your own with \emph{entfernen}
  \item \emph{Recall:} say the German for the dry cleaner, then the German for the launderette, then say \emph{entfernen} again
\end{itemize}

\begin{tcolorbox}[breakable,colback=teal!4,colframe=teal!35!black,title={Before you move on}]
What does entfernen mean? (\textquotedblleft{}To remove\textquotedblright{}.) Say it once more.
\end{tcolorbox}
\section[flicken]{flicken — to patch, to mend}

\label{lesson:GE-C446-flicken}



\begin{quote}
Before the new one: say the German for the launderette, then the German for to remove.
\end{quote}

\subsection*{You'll want to know: flicken}
\textbf{flicken} — \textquotedblleft{}to patch, to mend\textquotedblright{}.

\begin{quote}
Die Schneiderin flickt das Loch in der Jacke.
\end{quote}

\begin{quote}
\emph{The dressmaker patches the hole in the jacket.}
\end{quote}

\emph{Der Flicken} — the patch. \emph{Einen Fahrradschlauch flicken} — to mend a bike tube.

\begin{cousinweb}
\emph{flicken} is inherited from Middle High German \emph{vlicken} 'to put patches on' (EtymDB), built on \emph{vlëc} 'patch, spot', from Old High German \emph{flec} (EtymDB).
\end{cousinweb}

\begin{grammarlens}[title={What you've built}]
One more word for telling, explaining and describing at B1.
\end{grammarlens}

\subsection*{Your turn}
\begin{itemize}
\raggedright
  \item \emph{Say it:} \emph{flicken}
  \item \emph{Say it:} \emph{flicken}, once more
  \item \emph{Say it:} a sentence of your own with \emph{flicken}
  \item \emph{Recall:} say the German for the launderette, then the German for to remove, then say \emph{flicken} again
\end{itemize}

\begin{tcolorbox}[breakable,colback=teal!4,colframe=teal!35!black,title={Before you move on}]
What does flicken mean? (\textquotedblleft{}To patch, to mend\textquotedblright{}.) Say it once more.
\end{tcolorbox}
\section[trocknen]{trocknen — to dry}

\label{lesson:GE-C446-trocknen}



\begin{quote}
Before the new one: say the German for to remove, then the German for to patch.
\end{quote}

\subsection*{You'll want to know: trocknen}
\textbf{trocknen} — \textquotedblleft{}to dry\textquotedblright{}.

\begin{quote}
Im Sommer trocknet die Wäsche im Garten.
\end{quote}

\begin{quote}
\emph{In summer the washing dries in the garden.}
\end{quote}

Perfekt with \emph{sein} when something dries by itself: \emph{die Wäsche ist getrocknet}; with \emph{haben} when someone dries it.

\begin{cousinweb}
\emph{trocknen} is built on \emph{trocken} 'dry', inherited from Middle High German \emph{trucken} (EtymDB).
\end{cousinweb}

\begin{grammarlens}[title={What you've built}]
One more word for telling, explaining and describing at B1.
\end{grammarlens}

\subsection*{Your turn}
\begin{itemize}
\raggedright
  \item \emph{Say it:} \emph{trocknen}
  \item \emph{Say it:} \emph{trocknen}, once more
  \item \emph{Say it:} a sentence of your own with \emph{trocknen}
  \item \emph{Recall:} say the German for to remove, then the German for to patch, then say \emph{trocknen} again
\end{itemize}

\begin{tcolorbox}[breakable,colback=teal!4,colframe=teal!35!black,title={Before you move on}]
What does trocknen mean? (\textquotedblleft{}To dry\textquotedblright{}.) Say it once more.
\end{tcolorbox}
